\section{Unexecuted E5 replication and old-only engineering}\label{sec:e5-engineering}
A planned replication with three new Granite optimizer seeds and a disjoint
600-person same-corpus panel (E5) remains unexecuted: its strict readout gate
failed before search release. No E5 optimizer search, new-panel forecast or
holdout-target analysis was run. The engineering checks in
Table~\ref{tab:e5-engineering} use the original Cal120 forecasts and runtime
metadata only; they add neither a transfer sample nor a new optimizer result.
The reference remains the original calibrated-GA Granite seed-1
\texttt{base\_cal} archive (SHA-256 prefix \texttt{9ebeffc8f789}), with the
unchanged rule $\max|\Delta p|\leq5\times10^{-4}$.

\begin{table}[H]\centering\footnotesize
\caption{Old-only E5 engineering and independent readback. Completion or
verification of a diagnostic does not pass historical compatibility or release
search. Profile comparisons use their common respondent rows.}
\label{tab:e5-engineering}
\begin{tabular}{@{}p{.12\columnwidth}p{.82\columnwidth}@{}}\toprule
Job & Recorded result\\\midrule
23135 & Two fresh sequential servers, explicit batch-invariant OFF/ON:
30,720 confirmed forecasts and 24,258,640 prompt tokens.
OFF fails five of six profile pairs; ON probabilities are byte-identical
in all six.
Both arms fail all four comparisons with the canonical reference;
the maximum probability difference remains 0.0303456.\\
23637 & Independent CPU readback of archived raw top-20 log probabilities
and reconstruction of answer mass and normalized arrays, with checks of
prompt tokens, source lineage and costs, agrees with the diagnostic;
historical compatibility still fails.\\
23645 & All 42 historical Granite neutral-Cal arrays, two preserved replays
and eight OFF/ON profiles are retained. The canonical exact mode recurs in
six historical cells. The 42 historical full-Cal arrays have 27 exact modes
but one shared exact tail over respondents 9--120.\\
23654 & Fixed initialization-to-post-capture overlap across all 42 launches:
0/4 overlapping pairs and 27/857 nonoverlapping pairs share a mode;
proportion difference $-0.0315053$, two-sided task-label permutation $p=1$
(100,000 draws). The fixed overlap proxy has no positive support.\\\bottomrule
\end{tabular}\end{table}

The invariant ON path stabilizes this execution without recovering the
historical reference; it does not identify a single numerical kernel.
Repeated canonical bytes weaken a one-off corruption hypothesis without
proving absence or cause. The overlap test is observational, has only four
exposed pairs, and assumes task-label exchangeability that temporal or
launch-order dependence may violate. Its interval definition was not changed
after the result. Recorded selected attention backend, AOT action and graph
mode are constant across the 42 launches despite the multiple output modes.
An older fresh-server full-Cal first call and repeat with identical rendered
tokens (job 9149) already both fail the same canonical gate; this is not a new
E5 replication.
These checks establish no new-search-seed robustness or mechanism-level cause.
Full receipt hashes and claim locators are retained in
\texttt{E5\_ENGINEERING\_PAPER\_BINDING\_}\allowbreak
\texttt{20261024.json}.

The additional E5 engineering allocations (jobs 23036, 23121 and 23135)
consume 0.188333 GPU-hours, increasing the prior current-cycle total from
13.796111 to 13.984444 GPU-hours; job 23135 accounts for 0.094444 GPU-hours
of that increment. The independent CPU checks add no GPU allocation. A later mechanism-trace
allocation (24116) fails after one second, before server launch or forecasts,
adding 0.000278 GPU-hours and bringing the current-cycle total to
13.984722 GPU-hours. Its client-environment preflight validation requests
Triton package metadata unavailable in that interpreter; no hidden-state,
logit or sampler observation is obtained. That attempt is closed without a retry and supplies no layer-localization
evidence.
The separately inventoried historical minimum remains 300.599444 GPU-hours
(Table~\ref{tab:historical-cost}); it is outside this current-cycle total,
and the older job 9149 is outside the current-cycle ledger. Neither total is
an end-to-end project cost or an efficiency result.


A subsequent raw-only CPU check (24086) tests the narrower hypothesis of
unchanged digit-relative values plus a shared row offset and final
round-to-nearest FP32 storage. In the OFF full-Cal repeat contrast,
8 of 28,800 anchored digit components across four forecasts and two
respondents exceed the stated component-specific rounding bounds; the
maximum change is 0.250000119, versus a maximum bound of
$1.91\times10^{-6}$. All six within-ON contrasts are exactly zero.
This rejects that stated offset-and-final-rounding model, not every possible
internal arithmetic error, and identifies neither a kernel nor a canonical
repair. The failed JSON publication of its earlier implementation (24083)
is retained separately and is not interpreted as a numerical result.

Static opcode inspection, without deserializing or executing the AOT payload,
finds 21 embedded source texts in the current cached artifact at the
historically recorded path; all match the collected current sources
(15 vLLM and six Torch/Transformers files). This is a partial cached-source
witness for the compiled forward and support code. No October 7 artifact
hash exists, so it does not authenticate the historical bytes or the complete
scheduler, attention-kernel and runtime environment. Engineering-only
preflight failures (24108, 24110, 24112) and the successful CPU compile
check (24113) are retained; none is scientific parity evidence.


A separately authorized corrected trace (24315) validates Triton in the
pinned vLLM interpreter and the supervisor in its client interpreter
(CPU preflight 24313). One default-OFF server performs exactly two old-Cal120
reads: 14,400 confirmed forecasts and 11,369,760 prompt tokens. The GPU job
remains FAILED because its final join expects external request IDs, while
vLLM adds an eight-hex-character internal suffix. Its 14,400 captured rows
and raw artifacts are retained unchanged. A source-derived, posthoc CPU-only
binding correction (24328) requires unique raw and normalized IDs, an exact
14,400-row client/server bijection, recomputed request hashes, and unchanged
artifact/source hashes before and after analysis. The two additional source
files explaining the internal suffix are current post-run witnesses, absent
from the original seven-source runtime receipt. The earlier CPU binding
implementation failure (24322, before array analysis) is retained.

The bound trace contains 7,200 paired forecasts. Ten pairs differ already
in the exact BF16 LM-head-input hidden-state bits; all ten also differ
in returned top-logprob records. This localizes observed repeat
variation at or before that hidden-state boundary under instrumentation,
without identifying an earlier layer or kernel. All ten have different final
sampler-call shapes and repeat-local call ordinals, but these contexts also
differ widely among unchanged pairs and do not identify a causal batch effect.
The other 7,190 pairs have
identical hidden states and five captured canonical-digit logits; their
71,900 anchored digit-pair comparisons show no sampler-repeat variation
beyond the stated bound. No isolated head variation is observed with
identical hidden states. These nulls cover only selected coordinates and
the stated arithmetic model. Snapshot kernels and retained buffers can
alter scheduling; this trace does not establish the uninstrumented
historical cause, repair canonical compatibility, or release E5 search.
The one GPU allocation lasts 134 seconds (0.037222 GPU-hours), bringing
the current-cycle total to 14.021944 GPU-hours. No additional GPU replay or
new-target inference is performed for the CPU binding correction.


Retained-artifact CPU forensics (24574), prespecified over all 7,200 pairs,
finds zero token-signature, token-count or selected-position mismatches;
every selected position equals the prompt token count minus one. These are
receipt checks: the actual token-ID arrays were not retained. All ten changed
hidden-state pairs lie in the previously fixed first-eight-respondent group
(10/480 forecasts), versus 0/6,720 in the other 112 respondents. Each changed
vector differs numerically in 3,390--3,763 of 4,096 coordinates, with no
signed-zero-only changes; the least-squares common-scale residual is
0.7595--1.4154\% of the second vector's norm. These descriptive metrics do
not identify RMSNorm, KV-cache or batch scheduling as a cause. The tensors
correspond to selected final-RMSNorm outputs supplied as LM-head inputs;
the snapshots are queued only after the head call returns, and backend
nonmutation is not independently proved. Earlier-layer states were
not retained, so the earliest upstream divergence remains unobserved.


A CPU-only static AOT feasibility audit (24710; 203 seconds, two CPUs,
8 GiB, no GPU) inspects the unchanged 17,075,197-byte cached object using
pickle opcodes and Python ASTs without deserializing or executing payloads.
It inventories seven complete and five incomplete pickle units,
185 AST-readable Python-source records (34 proof-eligible), and 63 opaque
binary units, including
a 162,874-byte execution-code witness. The bounded audit does not prove complete producer--buffer--consumer maps
for the 66 boundaries, complete original numeric-launch association, or an
exact existing logits-index delivery channel; the runner's dynamic input expansion remains
unresolved. The verdict is FEASIBILITY\_HOLD / OBSERVER\_UNAVAILABLE, not a
proof that another static method could never recover these relationships.
The fixed protocol therefore stops before observer implementation or GPU
submission, with no eager fallback, signature change or numerical recompile.
